% ==================================================================
% id:       cond-005
% topic:    הסתברות מותנית ואי-תלות
% subtopic: הסתברות מותנית
% solution: full
% ==================================================================
נערך מחקר על הקשר בין עישון למחלה מסוימת. אלה הנתונים:
\begin{center}
\begin{tabular}{c|cc|c}
& חולה & לא חולה & סה״כ\\ \hline
מעשן & 25 & 75 & 100\\
לא מעשן & 20 & 180 & 200\\ \hline
סה״כ & 45 & 255 & 300
\end{tabular}
\end{center}
בוחרים באקראי אדם מתוך 300 הנבדקים.
\begin{enumerate}
\item מה ההסתברות שהוא מעשן?
\begin{solbox}
נסמן $S$ = „מעשן”, $D$ = „חולה”. כל הנבדקים סבירים באותה מידה, ולכן $\PP(S)=\frac{100}{300}=\ans{\frac13}$.
\end{solbox}
\item מה ההסתברות שהוא חולה?
\begin{solbox}
$\PP(D)=\frac{45}{300}=\ans{0.15}$.
\end{solbox}
\item מה ההסתברות שהוא מעשן וגם חולה?
\begin{solbox}
$\PP(S\cap D)=\frac{25}{300}=\ans{\frac1{12}}$.
\end{solbox}
\item ידוע שהנבדק מעשן. מה ההסתברות שהוא חולה?
\begin{solbox}
לפי ההגדרה ולפי סעיפים (א) ו-(ג):
\[
\PP(D\mid S)=\frac{\PP(D\cap S)}{\PP(S)}=\frac{1/12}{1/3}=\ans{\frac14}.
\]
אפשר גם לספור ישירות בתוך שורת המעשנים: מתוך 100 מעשנים, 25 חולים, כלומר $\frac{25}{100}$.
\end{solbox}
\end{enumerate}
